% ================================================================== Chapitre 7
\chapter{Thermodynamique de l'oxydoréduction}
\label{ch:redox}

Les potentiels d'électrode de la première année et l'enthalpie libre de réaction décrivent la même chose : la tendance d'un système rédox à évoluer. Ce chapitre fait le lien, ce qui permet de calculer des potentiels standard à partir de tables thermodynamiques, et inversement.

\section{Enthalpie libre d'une réaction d'oxydoréduction}

On considère la réaction entre le couple 1 ($\mathrm{Ox}_1/\mathrm{Red}_1$) et le couple 2 ($\mathrm{Ox}_2/\mathrm{Red}_2$), écrite pour échanger $\gamma$ électrons :
\[ (\mathcal{R})\qquad a_1\,\mathrm{Ox}_1 + b_2\,\mathrm{Red}_2 \;=\; b_1\,\mathrm{Red}_1 + a_2\,\mathrm{Ox}_2 \]

\begin{resultat}{Relation entre enthalpie libre de réaction et potentiels}
\[ \boxed{\Delta_rG = -\gamma F\,(E_1 - E_2)} \qquad\qquad \boxed{\Delta_rG^\circ = -\gamma F\,(E_1^\circ - E_2^\circ)} \]
$E_1$ et $E_2$ sont les potentiels de Nernst des deux couples calculés avec la composition \emph{actuelle} du système ; $F$ est la constante de Faraday.
\end{resultat}

\noindent\textbf{D'où vient le facteur $-\gamma F$ ?} Il est établi au chapitre~\ref{ch:piles} : lorsque la réaction se déroule dans une pile, l'avancement $\dd\xi$ fait circuler une charge $\gamma F\,\dd\xi$ sous la tension $E_1 - E_2$, et $\Delta_rG\,\dd\xi$ est égale au travail électrique reçu. On vérifie ici que la relation est \emph{cohérente} avec la loi de Nernst.

\begin{demo}[Cohérence avec la formule de Nernst]
En écrivant Nernst pour chaque couple (demi-équations multipliées pour faire intervenir $\gamma$ électrons) et en soustrayant :
\[ E_1 - E_2 = (E_1^\circ - E_2^\circ) + \frac{RT}{\gamma F}\ln\frac{a_{\mathrm{Ox}_1}^{a_1}\,a_{\mathrm{Red}_2}^{b_2}}{a_{\mathrm{Red}_1}^{b_1}\,a_{\mathrm{Ox}_2}^{a_2}} = (E_1^\circ - E_2^\circ) - \frac{RT}{\gamma F}\ln Q_r. \]
En multipliant par $-\gamma F$ : $-\gamma F(E_1-E_2) = -\gamma F(E_1^\circ-E_2^\circ) + RT\ln Q_r$. C'est exactement la structure $\Delta_rG = \Delta_rG^\circ + RT\ln Q_r$ : les deux relations encadrées sont compatibles entre elles et avec le chapitre~\ref{ch:evolution}.
\end{demo}

\begin{resultat}{Conséquences}
\begin{itemize}
  \item \textbf{Sens d'évolution} : $\Delta_rG < 0 \Leftrightarrow E_1 > E_2$. L'oxydant du couple de plus haut potentiel est réduit : c'est la justification thermodynamique de la règle du gamma.
  \item \textbf{Équilibre} : $\Delta_rG = 0 \Leftrightarrow E_1 = E_2$ : la solution n'a plus qu'un seul potentiel.
  \item \textbf{Constante d'équilibre} : $-RT\ln K^\circ = -\gamma F(E_1^\circ - E_2^\circ)$, d'où
  \[ \ln K^\circ = \frac{\gamma F}{RT}(E_1^\circ - E_2^\circ) \qquad\Longrightarrow\qquad \log K^\circ \underset{\SI{298}{K}}{=} \frac{\gamma(E_1^\circ - E_2^\circ)}{0{,}06}. \]
  \item Écriture compacte : $\Delta_rG = RT\ln\dfrac{Q_r}{K^\circ} = -\gamma F\,\Delta E$, avec $\Delta E = E_1 - E_2$.
\end{itemize}
\end{resultat}

\begin{ordres}
À $\SI{298}{K}$, un écart $\Delta E^\circ = \SI{0,1}{V}$ pour $\gamma = 1$ correspond à $\Delta_rG^\circ = \SI{-9,6}{kJ.mol^{-1}}$ et à $K^\circ\approx 10^{1{,}7}\approx 50$. Pour $\gamma = 2$ et $\Delta E^\circ = \SI{1,10}{V}$ (pile Daniell) : $\Delta_rG^\circ = \SI{-212}{kJ.mol^{-1}}$ et $K^\circ \approx 10^{37}$.
\end{ordres}

\section{Enthalpie libre standard d'une demi-équation}

\begin{definition}[enthalpie libre standard de demi-réaction]
À la demi-équation $\alpha\,\mathrm{Ox} + \gamma\,e^- + \dots = \beta\,\mathrm{Red} + \dots$, on associe par convention
\[ \boxed{\Delta_{1/2}G^\circ = -\gamma F\,E^\circ(\mathrm{Ox/Red})} \]
\end{definition}

Cette convention revient à mesurer chaque couple par rapport à l'ESH : $\Delta_{1/2}G^\circ$ est l'enthalpie libre standard de la réaction complète entre $\mathrm{Ox}$ et \ce{H2} (couple \ce{H+/H2}, de potentiel standard nul). Son intérêt :
\begin{itemize}
  \item pour toute réaction rédox, $\Delta_rG^\circ = \Delta_{1/2}G_1^\circ - \Delta_{1/2}G_2^\circ$ : les électrons « s'éliminent » comme dans l'équation ;
  \item les $\Delta_{1/2}G^\circ$ s'additionnent comme des grandeurs de réaction ordinaires, et l'on peut les combiner avec les $\Delta_rG^\circ$ d'autres réactions (dissolution, acido-basicité, formation).
\end{itemize}

\section{Calcul d'un potentiel standard inconnu}

\begin{methode}[potentiel standard d'un couple à partir de données]
\begin{enumerate}
  \item Écrire la demi-équation cherchée, avec ses $\gamma$ électrons.
  \item L'obtenir comme \emph{combinaison linéaire} de demi-équations connues et de réactions de constante connue ($K_s$, $K_a$, $K_e$, $\Delta_fG^\circ$\dots).
  \item Additionner les enthalpies libres standard : $-\gamma F E^\circ_{\text{cherché}} = \sum(\text{contributions})$, avec $\Delta_{1/2}G^\circ = -\gamma_kFE^\circ_k$ pour une demi-équation et $\Delta_rG^\circ = -RT\ln K^\circ$ pour une réaction.
  \item Diviser par $-\gamma F$.
\end{enumerate}
On n'additionne \textbf{jamais} les potentiels : ce sont les enthalpies libres qui s'additionnent.
\end{methode}

\subsection{Même nombre d'oxydation dans les deux couples}

\begin{exemple}[électrode au calomel : $E^\circ(\ce{Hg2Cl2}/\ce{Hg})$]
Données : \ce{Hg2^{2+} + 2e^- = 2 Hg(l)}, $E_1^\circ = \SI{0,797}{V}$ ; \ce{Hg2Cl2(s) = Hg2^{2+} + 2 Cl^-}, $\mathrm{p}K_s = 17{,}9$. Dans les deux couples, le mercure passe du n.o. $+\mathrm{I}$ à $0$ : seule l'espèce oxydée diffère (ion libre ou précipité).

\medskip
\noindent\emph{Méthode 1 -- égalité des potentiels.} Une solution au contact de \ce{Hg}, \ce{Hg2Cl2} et \ce{Cl-} n'a qu'un seul potentiel, donné par l'un ou l'autre couple :
\[ E = E_1^\circ + \frac{0{,}059}{2}\log\frac{[\ce{Hg2^{2+}}]}{c^\circ} \qquad\text{avec}\qquad \frac{[\ce{Hg2^{2+}}]}{c^\circ} = \frac{K_s}{([\ce{Cl-}]/c^\circ)^2} \]
\[ \Longrightarrow\quad E = \underbrace{E_1^\circ - \frac{0{,}059}{2}\,\mathrm{p}K_s}_{E^\circ(\ce{Hg2Cl2/Hg})} - 0{,}059\log\frac{[\ce{Cl-}]}{c^\circ}. \]
\emph{Méthode 2 -- combinaison des enthalpies libres.}
\[ \renewcommand{\arraystretch}{1.3}\begin{array}{cl@{\qquad}l}
 & \ce{Hg2Cl2(s) -> Hg2^{2+} + 2Cl-} & \Delta_rG^\circ = -RT\ln K_s\\
 + & \ce{Hg2^{2+} + 2e^- -> 2Hg(l)} & \Delta_{1/2}G^\circ = -2FE_1^\circ\\ \hline
 = & \ce{Hg2Cl2(s) + 2e^- -> 2Hg(l) + 2Cl-} & -2FE^\circ = -RT\ln K_s - 2FE_1^\circ
\end{array} \]
d'où $E^\circ = E_1^\circ + \dfrac{RT}{2F}\ln K_s = 0{,}797 - 0{,}0296\times17{,}9 = \SI{0,27}{V}$.
\end{exemple}

\begin{pieges}
Ici le terme $\frac{0{,}059}{2}\mathrm{p}K_s$ vaut $\SI{0,53}{V}$ : arrondir $0{,}059$ à $0{,}06$ introduit une erreur de $\SI{0,01}{V}$. Quand un coefficient est multiplié par un grand $\mathrm{p}K$, garder $0{,}059$.
\end{pieges}

\subsection{Nombres d'oxydation différents}

\begin{exemple}[$E^\circ(\ce{Fe^{3+}}/\ce{Fe})$]
Données : \ce{Fe^{3+} + e^- = Fe^{2+}}, $E_1^\circ = \SI{0,77}{V}$ ; \ce{Fe^{2+} + 2e^- = Fe(s)}, $E_2^\circ = \SI{-0,44}{V}$. La demi-équation cherchée \ce{Fe^{3+} + 3e^- = Fe(s)} est la somme des deux :
\[ -3F\,E_3^\circ = -1\cdot F\,E_1^\circ - 2F\,E_2^\circ \quad\Longrightarrow\quad E_3^\circ = \frac{E_1^\circ + 2E_2^\circ}{3} = \frac{0{,}77 - 0{,}88}{3} = \SI{-0,04}{V}. \]
\end{exemple}

\begin{figure}[H]
\centering
\begin{tikzpicture}[>=Stealth]
  \node[font=\large] (A) at (0,0) {\ce{Fe^{3+}}};
  \node[font=\large] (B) at (4,0) {\ce{Fe^{2+}}};
  \node[font=\large] (C) at (8,0) {\ce{Fe(s)}};
  \draw[->,line width=1.1pt,navy] (A) -- node[below,font=\small] {$\gamma=1$} node[above,font=\small] {$\SI{0,77}{V}$} (B);
  \draw[->,line width=1.1pt,navy] (B) -- node[below,font=\small] {$\gamma=2$} node[above,font=\small] {$\SI{-0,44}{V}$} (C);
  \draw[->,line width=1.1pt,c3] (A.north) to[out=35,in=145] node[above,font=\small] {$\gamma = 3$ : $E^\circ = \dfrac{1\times0{,}77 + 2\times(-0{,}44)}{3} = \SI{-0,04}{V}$} (C.north);
\end{tikzpicture}
\caption{Diagramme de potentiels du fer : les potentiels se combinent pondérés par le nombre d'électrons.}
\label{fig:latimer}
\end{figure}

\subsection{À partir des enthalpies libres de formation}

\begin{exemple}[$E^\circ(\ce{O2}/\ce{H2O})$]
Demi-équation : \ce{O2(g) + 4 H+ + 4 e^- = 2 H2O(l)}. Associée à l'ESH, elle correspond à la réaction \ce{O2(g) + 2 H2(g) -> 2 H2O(l)} échangeant $\gamma = 4$ électrons, avec $E^\circ(\ce{H+/H2}) = 0$ :
\[ \Delta_rG^\circ = 2\,\Delta_fG^\circ(\ce{H2O(l)}) = 2\times(-237{,}1) = \SI{-474,2}{kJ.mol^{-1}} = -4F\,(E^\circ - 0) \]
\[ \Longrightarrow\quad E^\circ(\ce{O2/H2O}) = \frac{474{,}2\times10^3}{4\times96485} = \SI{1,23}{V}. \]
\end{exemple}
